\chapter{失敗と混雑}
\label{ch:failures}

ヘッジファンドの歴史に残る失敗は、たいてい同じ形をしている。平時には安定して儲かる取引に、大きなレバレッジをかけ、同じ取引をしている人が多く、売りたいときに売れない。この章では、代表的な失敗を並べ、その共通の構造と、研究から分かる教訓をまとめる。

\section{共通の構造：解消のらせん}

\begin{figure}[H]
\centering
\begin{tikzpicture}[node distance=8mm and 12mm]
\node[box=c2,text width=30mm] (loss) {損失};
\node[box=c4,text width=30mm,right=24mm of loss] (margin) {担保の追加の要求・\\投資家の解約・\\リスクの上限};
\node[box=c7,text width=30mm,below=12mm of margin] (sell) {ポジションの\\投げ売り};
\node[box=c1,text width=30mm,below=12mm of loss] (price) {価格の下落\\（同じ取引をしている\\人も損をする）};
\draw[arr] (loss)--(margin);
\draw[arr] (margin)--(sell);
\draw[arr] (sell)--(price);
\draw[arr] (price)--(loss);
\node[lab,align=center] at ($(loss)!0.5!(sell)$) {レバレッジ・混雑・\\流動性の不一致が\\回転を速くする};
\end{tikzpicture}
\caption{解消のらせん（Brunnermeier・Pedersen 2009 \cite{brunnermeierpedersen2009}、Pedersen 2009 \cite{pedersen2009} の考え方の模式図）。}
\label{fig:spiral}
\end{figure}

\begin{itemize}
\item \textbf{レバレッジ}：小さな価格の変化で資本の大きな割合を失い、担保の追加を求められる。
\item \textbf{混雑}：同じ取引をしている人が多いと、一人の売りがほかの人の損失を生み、売りの連鎖が起きる。Stein（2009）\cite{stein2009} は、洗練された裁定者が、何人が同じ取引をしているかを知らないこと自体が、価格を不安定にすることを論じた。
\item \textbf{流動性の不一致}：持っている資産はすぐに売れないのに、資金（借入れ、投資家の資金）はすぐに引き上げられる。
\item \textbf{モデルの前提の外れ}：過去の相関・変動率・流動性を前提にしたリスクの物差し（VaR など）は、危機の局面では役に立たない。
\end{itemize}

Pedersen（2009）\cite{pedersen2009} は、これを「皆が同時に出口に走るとき」と呼び、2007〜2008 年の出来事をこの枠組みで説明した。

\section{LTCM（1998 年）}

LTCM は 1994 年に、Salomon Brothers の債券の裁定の部門の出身者が設立し、のちにノーベル経済学賞を受ける Merton と Scholes も加わった。主な取引は、債券の相対価値（第 \ref{ch:rv} 章）で、利回りの差が縮むことに賭けた。

\begin{itemize}
\item 1998 年の初めに、資本の 25 倍前後の資産を持ち、デリバティブの想定元本はそれよりはるかに大きかった。
\item 1998 年 8 月のロシアの債務不履行の後、投資家が一斉に安全で流動性の高い資産（最新の米国債など）に逃げ、LTCM が賭けていたほぼすべての差が同時に広がった。
\item ほかの多くの銀行・ヘッジファンドが似たポジションを持っていて、LTCM の窮状が知られると、その取引の反対側に立つ人が現れた。
\item 9 月、ニューヨーク連邦準備銀行の仲介で、大手の金融機関が資本を注入し、ファンドは管理の下で清算された。
\end{itemize}

Jorion（2000）\cite{jorion2000} は、LTCM のリスク管理を分析し、短い期間の平穏なデータで推定した変動率と相関に基づく VaR が、危機の局面のリスクを大きく過小評価したこと、レバレッジが小さな価格の変化を致命的にしたこと、流動性のリスクと、ほかの参加者の行動のリスクが測られていなかったことを指摘した。Edwards（1999）\cite{edwards1999} は、この出来事が金融システムに与えた影響と、規制の論点を整理した。

\section{Amaranth（2006 年）}

Amaranth Advisors は、マルチストラテジーのファンドだったが、一人のトレーダーの天然ガスの先物の取引（冬と春の限月の価格差に賭けるスプレッドの取引）が、ファンドの資産の大部分のリスクを占めるようになっていた。2006 年 9 月、価格差が逆に動き、数週間で 60 億ドルを超える損失を出して、ファンドは閉じた。

Chincarini（2007）\cite{chincarini2007} と Till（2007）\cite{till2007} は、ポジションが市場全体の建玉のかなりの割合を占めるほど大きく、損を限るために売ろうとしても、自分の売りが価格をさらに不利に動かす状態にあったこと、リスクの物差しがこの流動性の問題を捉えていなかったことを分析した。一つのファンドの中でも、一つの取引への集中と、市場の厚さを超えるポジションが、致命的になりうる。

\section{2007 年 8 月と、2008 年}

2007 年 8 月のクオンツの急落（第 \ref{ch:statarb} 章、Khandani・Lo 2011 \cite{khandanilo2011}）は、市場全体が平穏な中で、混雑したポジションの解消の連鎖が起きた例である。

2008 年の金融危機では、ヘッジファンドの業界の全体が大きく損をした（業界の指数で年 2 割前後の損失と広く報じられた）。

\begin{itemize}
\item \textbf{Ben-David・Franzoni・Moussawi（2012）}\cite{bfm2012}：2008 年の後半に、ヘッジファンドは株の保有を大きく減らした（約 3 割）。売りは、投資家の解約と、資金の貸し手からの圧力によるもので、ほかの機関投資家より大きかった。
\item \textbf{Aragon・Strahan（2012）}\cite{aragonstrahan2012}：Lehman Brothers をプライムブローカーにしていたファンドは、資産が凍結されて売買できなくなり、破綻しやすくなった。そうしたファンドが持っていた株は、流動性が大きく下がった。ヘッジファンドが平時に流動性の供給者として果たしている役割を、逆側から示した。
\item \textbf{Mitchell・Pulvino（2012）}\cite{mitchellpulvino2012}：プライムブローカーが資金の提供を引き上げたことが、転換社債の裁定や合併裁定の価格の崩れを引き起こした（第 \ref{ch:rv} 章）。
\end{itemize}

大手のマルチストラテジーのファンドの中にも、2008 年に資産の半分前後を失ったものがあった、と報じられている。

\section{ボラティリティの売り（2018 年 2 月）と国債（2020 年 3 月）}

2018 年 2 月の VIX の急騰による、ボラティリティの売りの上場商品の崩壊（Augustin ほか 2021 \cite{volmageddon2021}）と、2020 年 3 月の国債の市場の混乱でのヘッジファンドの投げ売り（Kruttli ほか 2021 \cite{kruttli2021}）は、第 \ref{ch:rv} 章で扱った。どちらも、仕組み（日々のリバランス、担保の条件）が、価格の変化を増幅する方向の売買を機械的に生んだ。

\section{クオンツの冬（2018〜2020 年）}

割安を中心とする多くのファクターが同時に長く損をした（Blitz 2021 \cite{blitz2021}、第 \ref{ch:equityfactor} 章）。数日の急落ではなく、数年にわたる低迷で、多くのクオンツのファンドが資金を失った。リスク・プレミアムの戦略でも、何年も損が続くことがあり、そのあいだ投資家が資金を預け続けるとは限らない。

\section{2021 年：空売りの踏み上げと、隠れたレバレッジ}

\begin{itemize}
\item \textbf{GameStop}：2021 年 1 月、インターネットの掲示板で集まった個人投資家が、空売りの多い GameStop の株を買い上げ、株価は数週間で何十倍にもなった。大きな空売りを持っていたヘッジファンドの Melvin Capital は、1 月だけで資産の半分を超える損失を出し、ほかのファンドから資本の注入を受けたと報じられた。空売りには、損失に上限がないという固有のリスクがあり、空売りの残高そのものが「混雑」の指標になる。
\item \textbf{Archegos}：2021 年 3 月、ファミリー・オフィス（一族の資産の運用会社）の Archegos Capital Management が、複数の銀行と結んだトータル・リターン・スワップ（株を持たずに、株の値動きを受け取る契約）で、少数の銘柄に巨大なレバレッジのポジションを持っていた。株価が下がって担保の追加に応じられなくなると、銀行が一斉に担保の株を売り、銀行の損失は合わせて 100 億ドルを超えた（Credit Suisse だけで約 55 億ドルと報じられた）。スワップによるポジションは開示されず、各銀行は、ほかの銀行が同じ顧客に同じポジションを持たせていることを知らなかった。
\end{itemize}

これらは、学術研究の分析がまだ少ない最近の出来事で、事実関係は報道と当局の報告に基づく。

\section{混雑の測り方}

混雑は直接は見えないが、いくつかの測り方が提案されている。

\begin{itemize}
\item \textbf{保有に基づく}：機関投資家の保有の開示（米国の 13F）、空売りの残高、ファンドの資金の流れから、同じ銘柄・同じファクターに集まっている資金を測る。
\item \textbf{リターンに基づく}：Lou・Polk（2022）\cite{loupolk2022} は、モメンタムの勝ち組（あるいは負け組）の銘柄どうしの、ファクターで説明できないリターンの相関（コモメンタム）が、モメンタムの裁定者の多さを表すとした。コモメンタムが高い時期の後には、モメンタムの収益が低く、のちに反転が起きる。裁定者が多すぎると、価格は行き過ぎる。
\item \textbf{価格に基づく}：ファクターの上位と下位の銘柄の評価の差（割安の差の広がり）。混雑すると評価の差が縮む。
\end{itemize}

\section{効果の減衰}

公表の後の減衰（McLean・Pontiff 2016 \cite{mcleanpontiff2016}、第 \ref{ch:equityfactor} 章）は、「ゆっくりした混雑」である。Falck・Rej・Thesmar（2021）\cite{falck2022} は、公表されたアノマリーの標本外のシャープレシオの低下を、事前に分かる特徴で予測した。公表の年だけで低下の分散の 30\% を説明し、新しく公表されるファクターほど低下が大きい（毎年 5 ポイントずつ大きくなる）。シグナルの計算に要する操作の数や、標本の中のシャープレシオの外れ値への敏感さ（過学習に関係する変数）が、さらに 15\% を説明した。裁定に関係する変数の説明力は小さかった。

\section{教訓}

\begin{keybox}[失敗の研究から分かること]
\begin{enumerate}
\item 平時のシャープレシオは、まれな大損のリスクを示さない。最悪の場合の損失と、それに耐えられる資金の構造（レバレッジ、資金の引き上げの条件）を見る。
\item 自分の戦略が正しくても、同じ取引をしているほかの人の事情で損をする。有名な論文のシグナル、広く使われる道具（LLM の語調を含む）ほど、混雑しやすい。
\item 持っている資産の流動性と、資金の流動性を合わせる。
\item 過去のデータで推定したリスクの物差しは、危機の局面で外れる。ストレス・テストで補う。
\item 公表された戦略の効果は、時間とともに小さくなる。新しいものほど、早く小さくなる。
\end{enumerate}
\end{keybox}
